The quadratic family of Section~\ref{sec:polynomial-threaded-carriers}
shows that one polynomial can define both a real carrier and a finite complex
thread.  Its two-strand monodromy is nevertheless abelian.  This section studies
the first example in which the same mechanism has genuinely noncommutative
range.  The point is not the degree six by itself.  It is the coexistence of
three polynomial layers:
\begin{equation}
 P(x),
 \qquad
 P(x)-t,
 \qquad
 Q_P(t):=6^{-6}\operatorname{disc}_x(P(x)-t).
 \label{eq:sextic-three-polynomial-layers}
\end{equation}
The first is a spatial polynomial, the second relates its slices, and the third
records exactly where those slices become singular.  Passing from $P$ to $Q_P$
extracts the event schedule while forgetting how the six sheets are attached.

All Lyashko--Looijenga, braid, Birman--Hilden, Picard--Lefschetz, and
Galois-theoretic inputs in this section are classical.  The purpose here is to
put them in one auditable AEG example and to identify precisely which
information each passage retains or forgets.  No claim is made that AEG
created those classical monodromy theorems.

\subsection{The event polynomial and a 1296-sheet forgetting map}

Let
\begin{align}
 \mathcal P_6^{\mathrm{mc}}
 &=\left\{
 x^6+a_2x^4+a_3x^3+a_4x^2+a_5x+a_6
 \right\}\cong\C^5,
 \label{eq:monic-centered-sextic-space}\\
 \mathcal Q_5
 &=\left\{
 t^5+b_1t^4+b_2t^3+b_3t^2+b_4t+b_5
 \right\}\cong\C^5.
 \label{eq:monic-quintic-event-space}
\end{align}
The superscript $\mathrm{mc}$ means monic with vanishing $x^5$ coefficient;
it does not remove the constant term.  For $P\in\mathcal P_6^{\mathrm{mc}}$,
define
\begin{equation}
 \operatorname{LL}_6(P)=Q_P(t)
 :=6^{-6}\operatorname{disc}_x(P(x)-t).
 \label{eq:sextic-ll-map}
\end{equation}
If $c_1,\ldots,c_5$ are the critical points of $P$, counted with
multiplicity, then the resultant identity gives
\begin{equation}
 Q_P(t)=\prod_{i=1}^5\bigl(t-P(c_i)\bigr).
 \label{eq:event-polynomial-critical-values}
\end{equation}
Thus $Q_P$ is the monic polynomial of critical values.  Let
\begin{equation}
 \mathcal U_5
 =\{Q\in\mathcal Q_5:\operatorname{disc}_tQ\ne0\},
 \qquad
 \mathcal P_6^\circ=\operatorname{LL}_6^{-1}(\mathcal U_5).
 \label{eq:ll-regular-loci}
\end{equation}

The completely explicit basepoint is
\begin{equation}
 P_0(x)=x^6-x.
 \label{eq:sextic-base-polynomial}
\end{equation}
Write $\zeta=e^{2\pi i/5}$.

\begin{theorem}[The sextic event schedule and its LL fiber]
\label{thm:sextic-ll-event-schedule}
The following statements hold.
\begin{enumerate}
\item The critical points and critical values of $P_0$ are
  \begin{equation}
   c_k=6^{-1/5}\zeta^k,
   \qquad
   v_k=P_0(c_k)=-5\,6^{-6/5}\zeta^k,
   \qquad k\in\Z/5\Z.
   \label{eq:sextic-critical-points-values}
  \end{equation}
  They are simple and pairwise distinct, and
  \begin{equation}
   \operatorname{disc}_x(x^6-x-t)=6^6t^5+5^5,
   \qquad
   Q_0(t)=t^5+\frac{5^5}{6^6}.
   \label{eq:sextic-explicit-discriminant}
  \end{equation}
\item The restriction
  \begin{equation}
   \operatorname{LL}_6:\mathcal P_6^\circ\longrightarrow\mathcal U_5
   \label{eq:sextic-ll-etale-cover}
  \end{equation}
  is a finite \`etale covering of degree
  \begin{equation}
   6^{6-2}=6^4=1296.
   \label{eq:sextic-ll-degree}
  \end{equation}
  Consequently the particular fiber
  $\operatorname{LL}_6^{-1}(Q_0)$ consists of exactly $1296$ reduced points.
\item The group $\mu_6$ acts freely on $\mathcal P_6^\circ$ over
  $\mathcal U_5$ by
  \begin{equation}
   (\xi\cdot P)(x)=P(\xi x),
   \qquad \xi^6=1.
   \label{eq:sextic-source-rotation-action}
  \end{equation}
  Hence the $1296$ coordinate-normalized lifts form $216$ orbits under the
  residual source-coordinate rotations.  This is the first quotient by the
  allowed source reparametrizations at fixed target; it is still not a count
  of pairwise nonisomorphic genus-two curves.  In particular, $1296$ is not a
  count of pairwise nonisomorphic curves or polynomial covers.
\item After choosing a base fiber labeling and distinguished paths to the five
  critical values, every LL lift determines a reduced factorization of a fixed
  six-cycle into five transpositions.  There are $1296$ such factorizations;
  the five transposition edges form a spanning tree on six labeled vertices.
  This is a marked combinatorial description, not a canonical unlabeled tree
  attached to $P$.
\end{enumerate}
\end{theorem}

\begin{proof}
The equation $P_0'(c)=6c^5-1=0$ gives $c^5=1/6$, and then
\begin{equation}
 P_0(c)=c^6-c=\frac{c}{6}-c=-\frac56c.
 \label{eq:sextic-critical-value-computation}
\end{equation}
Thus $v_k^5=-5^5/6^6$, proving the first formula and
\eqref{eq:sextic-explicit-discriminant}.

The finite covering statement and the degree $n^{n-2}$ for monic centered
degree-$n$ polynomials are the classical Lyashko--Looijenga theorem
\cite{Looijenga1974Bifurcation,Lyashko1984Bifurcation}.  For completeness, the
degree is also visible from the weighted coordinates.  If $Q_P=t^5$, all
critical values vanish; this forces $P=(x-r)^6$, and the missing $x^5$ term
forces $r=0$.  Thus the quasihomogeneous LL map is finite.  The source coefficients
$a_2,\ldots,a_6$ have weights $2,3,4,5,6$, whereas the five target coefficients
have weights $6,12,18,24,30$.  Their product ratio is
\begin{equation}
 \frac{6\cdot12\cdot18\cdot24\cdot30}
      {2\cdot3\cdot4\cdot5\cdot6}=1296.
 \label{eq:sextic-ll-weighted-degree}
\end{equation}
At a polynomial with five distinct critical points, the differential of the
ordered critical-value map sends a variation $\dot P\in\C[x]_{\le4}$ to
$(\dot P(c_1),\ldots,\dot P(c_5))$.  This is a Vandermonde isomorphism, so the
map is \`etale over $\mathcal U_5$.  Since $Q_0$ has five distinct roots, its
fiber has the stated reduced cardinality.

Equation \eqref{eq:sextic-source-rotation-action} preserves $Q_P$.  If a
nontrivial $\xi$ fixes $P$, then $P$ is a polynomial in $x^r$ for
$r=\operatorname{ord}(\xi)>1$.  Either zero is a degenerate critical point or
distinct critical points share a critical value; hence such a $P$ does not lie
in $\mathcal P_6^\circ$.  The action is therefore free there.  Finally, the
factorization count is D\'enes' theorem
\cite{Denes1959TranspositionFactorizations}; each factorization has a spanning
tree as its edge graph.  The auxiliary ordering and labeling explain the last
qualification.
\end{proof}

For $P_0$ itself the two pentagons have a simple geometric relation.  The six
roots at $t=0$ are a central root and the fifth roots of unity, while each
critical point lies on one of the five radial segments joining them.  The
critical values form a smaller regular pentagon in the $t$-plane, rotated by
$\pi$.

\begin{figure}[ht]
\centering
\begin{tikzpicture}[scale=1.05,>=Stealth]
  \begin{scope}[shift={(-3.15,0)}]
    \node[font=\small\bfseries] at (0,2.0) {$x$-plane: roots and vanishing tree};
    \fill (0,0) circle (2pt) node[below left] {$0$};
    \foreach \k in {0,...,4} {
      \coordinate (b\k) at ({1.55*cos(72*\k)},{1.55*sin(72*\k)});
      \coordinate (c\k) at ({1.08*cos(72*\k)},{1.08*sin(72*\k)});
      \draw[blue!55!black,thick] (0,0)--(b\k);
      \fill (b\k) circle (2pt);
      \fill[red!70!black] (c\k) circle (1.7pt);
    }
    \node[blue!55!black,font=\scriptsize] at (0,-1.93)
      {five radial vanishing arcs};
    \node[red!70!black,font=\scriptsize] at (0,-2.2)
      {red points: $c_k$};
  \end{scope}
  \begin{scope}[shift={(3.15,0)}]
    \node[font=\small\bfseries] at (0,2.0) {$t$-plane: event polynomial $Q_0$};
    \draw[gray!55] (0,0) circle (1.25);
    \foreach \k in {0,...,4} {
      \coordinate (v\k) at ({1.25*cos(72*\k+180)},{1.25*sin(72*\k+180)});
      \fill[red!70!black] (v\k) circle (2pt);
      \draw[->,gray!55] (0,0)--(v\k);
    }
    \fill (0,0) circle (1.5pt) node[below left] {$t_*=0$};
    \node[red!70!black,font=\scriptsize] at (0,-1.93)
      {$Q_0(t)=0$ at the five $v_k$};
  \end{scope}
\end{tikzpicture}
\caption{The explicit sextic calibration.  The left star records which two
roots collide at each event; the right pentagon records when the events occur.
The LL map retains the right-hand schedule and forgets the left-hand gluing
data.}
\label{fig:sextic-star-and-events}
\end{figure}

\subsection{A nested discriminant records events of events}

The first discriminant produces $Q_P$.  Its value and its own discriminant
then extract two further layers:
\begin{equation}
 Q_P(0)=6^{-6}\operatorname{disc}_xP,
 \qquad
 \Delta_{\mathrm{LL}}(P):=\operatorname{disc}_tQ_P.
 \label{eq:sextic-nested-discriminants}
\end{equation}
The first equality detects whether the zero slice $P(x)=0$ has colliding roots.
The second detects whether the event schedule itself has collided.

Let $\mathcal C$ be the reduced caustic hypersurface on which $P'$ has a
multiple root, and let $\mathcal M$ be the reduced Maxwell hypersurface on
which two distinct critical points have the same critical value.  Choose
reduced equations $C_6$ and $M_6$ for them.

\begin{proposition}[The double-discriminant divisor]
\label{prop:sextic-double-discriminant}
There is a nonzero constant $\kappa$ such that
\begin{equation}
 \operatorname{disc}_tQ_P
 =\kappa\,C_6(P)^3M_6(P)^2.
 \label{eq:sextic-maxwell-caustic-factorization}
\end{equation}
Equivalently, the first divisor equality below holds.  Moreover, the
ramification divisor is
\begin{equation}
 \operatorname{LL}_6^*(\text{target discriminant})
 =3\mathcal C+2\mathcal M,
 \qquad
 R_{\operatorname{LL}_6}=2\mathcal C+\mathcal M.
 \label{eq:sextic-ll-divisor-multiplicities}
\end{equation}
At a generic Maxwell point the pencil $y^2=P(x)-t$ has two nodes in one
fiber.  At a generic caustic point it has a cusp modeled by $y^2=z^3$.
For $P_0$, neither event occurs: its five singular spectral fibers each have
exactly one node.
\end{proposition}

\begin{proof}
Set-theoretically, two roots of $Q_P$ collide precisely when two distinct
critical points share a value or when a critical point degenerates.  At a
generic Maxwell point, the critical-value difference is a local parameter
$u$, while the unordered target collision coordinate is its square $u^2$.
At a generic caustic point the local polynomial is $z^3-uz$; its two critical
values differ by a nonzero multiple of $u^{3/2}$, so their squared difference
has order three.  These local models give the multiplicities two and three in
\eqref{eq:sextic-maxwell-caustic-factorization}, and subtracting one gives the
ramification multiplicities in \eqref{eq:sextic-ll-divisor-multiplicities}.
This is the classical LL bifurcation calculation
\cite{Looijenga1974Bifurcation,Lyashko1984Bifurcation}.

For the pencil $y^2=P(x)-t$, a repeated critical value gives two distinct
double roots at the same $t$, hence two nodes.  A double critical point gives
the cubic local equation.  The explicit critical points in
\eqref{eq:sextic-critical-points-values} are simple and their values are
distinct, proving the last assertion.
\end{proof}

The scalar $\Delta_{\mathrm{LL}}$ forgets even more than $Q_P$: all $1296$
points in an LL fiber have exactly the same $Q_P$ and the same second
discriminant.  This makes the tower useful as an information filtration, but
it cannot distinguish the realizations that it has forgotten.

\subsection{The star generates every six-braid}

At the regular parameter $t_*=0$ the roots of $P_0(x)-t_*$ are
\begin{equation}
 b_0=0,
 \qquad
 b_{k+1}=\zeta^k,
 \qquad k=0,\ldots,4.
 \label{eq:sextic-base-root-labels}
\end{equation}
Let $\gamma_k$ be the radial path from $0$ to $v_k$ in the $t$-plane, and let
\begin{equation}
 \alpha_k=\{r\zeta^k:0\le r\le1\}
 \label{eq:sextic-star-vanishing-arcs}
\end{equation}
be the corresponding radial arc in the $x$-plane.

\begin{theorem}[Full braid monodromy of one sextic pencil]
\label{thm:sextic-full-braid-monodromy}
Let $V=\{v_0,\ldots,v_4\}$.  The root monodromy of the single pencil
$P_0(x)-t$ is a surjection
\begin{equation}
 \mu_{P_0}:\pi_1(\C\setminus V,0)\cong F_5
 \twoheadrightarrow B_6.
 \label{eq:sextic-root-monodromy-surjection}
\end{equation}
More precisely, a positive meridian about $v_k$ maps to the positive
half-twist $h_k$ along $\alpha_k$, with permutation shadow
$(0,k+1)$.  The five half-twists $h_0,\ldots,h_4$ generate $B_6$.
\end{theorem}

\begin{proof}
Along the $k$th radial ray one has
\begin{equation}
 P_0(r\zeta^k)=\zeta^k(r^6-r).
 \label{eq:sextic-radial-restriction}
\end{equation}
As $t$ moves from $0$ to $v_k$, the roots beginning at $0$ and $\zeta^k$
therefore move along $\alpha_k$ and collide at $c_k$.  The critical point is
simple, so Picard--Lefschetz local monodromy is the positive half-twist along
that arc.  The five arcs form an embedded spanning star on the six roots.
Half-twists along the edges of an embedded spanning tree generate the braid
group: conjugating one star half-twist by another produces a half-twist between
two outer vertices, from which a standard Artin chain is obtained
\cite[Chapter~1]{Birman1974Braids}.  Hence the image is all of $B_6$.
\end{proof}

This is the first nonabelian universality statement in the paper that is
generated by one fixed displayed polynomial rather than by an arbitrary
supplied polynomial family.  Every six-braid occurs as the root transport of
some loop in the five-punctured $t$-plane.  Closing that braid can realize every
link of braid index at most six, but closure and Markov stabilization remain
external topological operations; the theorem itself is a statement about
embedded root transport.

\subsection{One loop, four compatible geometric readings}

Let $\gamma:S^1\to\C\setminus V$ be a smooth loop.  The same expression
\begin{equation}
 \mathcal P_\gamma(x,s)=P_0(x)-\gamma(s)
 \label{eq:sextic-loop-polynomial}
\end{equation}
has four related zero objects:
\begin{align}
 A_\gamma(x,s)&=\im\mathcal P_\gamma(x,s),
 &\Sigma_\gamma&=A_\gamma^{-1}(0),
 \label{eq:sextic-real-carrier}\\
 K_\gamma&=\mathcal P_\gamma^{-1}(0)\subset\Sigma_\gamma,
 &C_s&:\ y^2=\mathcal P_\gamma(x,s).
 \label{eq:sextic-thread-spectral-curve}
\end{align}

\begin{proposition}[Carrier, thread, and spectral mapping torus]
\label{prop:sextic-carrier-thread-spectral}
Assume that $\gamma$ meets each horizontal wall
\begin{equation}
 \{t\in\C:\im t=\im v_k\}
 \label{eq:sextic-real-critical-walls}
\end{equation}
transversely whenever it meets it.  Then:
\begin{enumerate}
\item $K_\gamma\to S^1$ is a square-free six-sheet thread contained in
  $\Sigma_\gamma$, and its braid is $\mu_{P_0}([\gamma])$.
\item The projection $\Sigma_\gamma\to S^1$ is critical exactly at
  $(c_k,s)$ with $\im(v_k-\gamma(s))=0$.  Each such point is a saddle
  reconnection.  These real carrier events are wall crossings, not collisions
  of the six roots.
\item On the complement of the five spatial critical points, every slice is a
  singular AES with
  \begin{equation}
   g_s=
   \frac{|6x^5-1|^2}
        {\mu^2+\lambda^2A_\gamma(x,s)^2}|\dd x|^2.
   \label{eq:sextic-aeg-slice-metric}
  \end{equation}
\item Let $D\subset\C$ be a disc containing the six roots throughout the
  loop.  The double cover of $D\times S^1$ branched over $K_\gamma$ is the
  mapping torus of the Birman--Hilden lift of
  $\mu_{P_0}([\gamma])$ to the genus-two surface with two boundary
  components.  Capping the two boundary components gives the mapping torus of
  the compact genus-two spectral monodromy.  This filling is along the two
  fiber-boundary slopes; downstairs it fills $D\times S^1$ to
  $S^2\times S^1$, not to the ordinary $S^3$ braid closure.
\end{enumerate}
\end{proposition}

\begin{proof}
Because $\gamma$ avoids $V$, every root of $P_0(x)-\gamma(s)$ is simple;
Proposition~\ref{prop:polynomial-carrier-thread} gives the finite thread and
identifies the vertical critical points of the carrier.  A simple holomorphic
critical point has local imaginary part equivalent to $\im(z^2)$, so a
transverse parameter wall crossing gives a Morse saddle.  The metric is the
specialization of \eqref{eq:polynomial-carrier-aeg-metric}.

The double cover of a disc branched at six points has Euler characteristic
$2\chi(D)-6=-4$ and two boundary components, hence it is $\Sigma_{2,2}$.
Transporting this cover around $\gamma$ gives exactly the mapping torus of the
lifted root braid.  Capping the two boundary circles compactifies
$y^2=P_0(x)-\gamma(s)$ by its two points over infinity and fills the two
boundary tori by $D^2\times S^1$ along the fiber-boundary slopes.
\end{proof}

The distinction in item~(2) is essential.  The spectral curve becomes nodal
only at the point event $\gamma(s)=v_k$.  A loop avoiding $v_k$ can nevertheless
cross its horizontal wall twice, producing two real carrier saddles, and can
wind around $v_k$, producing nontrivial root and spectral monodromy.  The same
event polynomial controls all three phenomena, but they are different strata.

\subsection{Genus two sees the full mapping class and symplectic groups}

Let $q:\Sigma_{2,2}\to D$ be the double cover branched at the six roots in
\eqref{eq:sextic-base-root-labels}.  The lift of $\alpha_k$ is a simple closed
curve $d_k=q^{-1}(\alpha_k)$, and the half-twist $h_k$ lifts to the right Dehn
twist $T_{d_k}$.

\begin{theorem}[Full genus-two and Gauss--Manin monodromy]
\label{thm:sextic-full-genus-two-monodromy}
The sextic pencil has the following monodromy filtration.
\begin{equation}
 F_5\twoheadrightarrow B_6
 \hookrightarrow\operatorname{Mod}(\Sigma_{2,2};\partial)
 \longrightarrow\operatorname{Mod}(\Sigma_2)
 \longrightarrow\operatorname{Sp}_4(\Z)
 \longrightarrow\operatorname{Sp}_4(\F_2)\cong S_6.
 \label{eq:sextic-monodromy-filtration}
\end{equation}
The first arrow is Theorem~\ref{thm:sextic-full-braid-monodromy}; the second is
the Birman--Hilden injection.  After capping the boundary, the composite
\begin{equation}
 F_5\twoheadrightarrow\operatorname{Mod}(\Sigma_2)
 \label{eq:sextic-full-mapping-class-surjection}
\end{equation}
is surjective.  The integral Gauss--Manin image of the pencil is all of
$\operatorname{Sp}_4(\Z)$.  Modulo two, the symplectic action is exactly the
root permutation, so there is an exact sequence
\begin{equation}
 1\longrightarrow PB_6\longrightarrow B_6
 \longrightarrow\operatorname{Sp}_4(\F_2)\cong S_6
 \longrightarrow1.
 \label{eq:sextic-mod-two-forgetting-sequence}
\end{equation}
The compactified family first retains the two distinguished points
$p_+,p_-$ over infinity.  Passing to the displayed unmarked mapping class group
forgets them.  The kernel of the closed-surface map from $B_6$ is strictly
larger than its center; in particular no isomorphism
$B_6/Z(B_6)\cong\operatorname{Mod}(\Sigma_2)$ is asserted.
\end{theorem}

\begin{proof}
The Birman--Hilden lifting theorem gives the injection for a double branched
cover of the disc and sends half-twists to Dehn twists
\cite{BirmanHilden1973Isotopies}.  The star half-twists generate all of $B_6$,
so their image contains the lifts of a standard Artin chain.  After capping,
those standard chain twists generate the full genus-two mapping class group;
equivalently, every genus-two mapping class is hyperelliptic.  The classical
integral symplectic representation of these twists is surjective onto
$\operatorname{Sp}_4(\Z)$
\cite{ACampo1979Tresses}.  Finally, the standard model
\begin{equation}
 H_1(\Sigma_2;\F_2)
 \cong
 \left\{(z_0,\ldots,z_5)\in\F_2^6:\sum z_i=0\right\}
 \big/\langle(1,\ldots,1)\rangle
 \label{eq:genus-two-mod-two-root-model}
\end{equation}
identifies a half-twist with the transvection induced by the corresponding
root transposition.  This realizes the exceptional isomorphism
$\operatorname{Sp}_4(\F_2)\cong S_6$ and proves
\eqref{eq:sextic-mod-two-forgetting-sequence}.
\end{proof}

Thus one sparse integral pencil realizes every genus-two mapping class, and
therefore periodic, reducible, and pseudo-Anosov monodromies all occur.  The
smooth fibers are still all homeomorphic to one genus-two surface.  The
diversity lies in gluing histories and their three-dimensional mapping tori,
not in a changing fiber genus.

As a concrete dynamical witness, the braid
\begin{equation}
 \beta=\sigma_1\sigma_3\sigma_5\sigma_2^{-1}\sigma_4^{-3}
 \label{eq:sextic-penner-witness-braid}
\end{equation}
is realized by some loop in $\C\setminus V$.  Its lift is the product of
positive twists on the odd curves of the standard five-chain and negative
twists on the even curves, with every component represented.  Penner's
criterion therefore makes it pseudo-Anosov
\cite{Penner1988PseudoAnosov}.  In the standard symplectic representation its
matrix, in the Digne--Kassel symplectic basis with column-vector convention
\cite{DigneKassel2023SymplecticSteinberg}, is
\begin{equation}
 M_\beta=
 \begin{pmatrix}
 3&-3&2&-1\\
 -1&7&-1&2\\
 1&0&1&0\\
 0&3&0&1
 \end{pmatrix}
 \in\operatorname{Sp}_4(\Z)
 \label{eq:sextic-penner-homological-matrix}
\end{equation}
Its characteristic polynomial is
\begin{equation}
 X^4-12X^3+31X^2-12X+1,
 \label{eq:sextic-penner-homological-polynomial}
\end{equation}
which is irreducible over $\Q$: after division by $X^2$ and the substitution
$Y=X+X^{-1}$ it becomes $Y^2-12Y+29$, whose discriminant is $28$.
The alternative product of two monic quadratics with constant terms $-1$
has opposite $X^3$ and $X$ coefficients, and $X=\pm1$ is not a root, so no
other rational factorization occurs.
Reversing braid or action conventions replaces the matrix by an inverse,
transpose, or conjugate without changing this reciprocal characteristic
polynomial.  This polynomial is a homological shadow; its
spectral radius is only a lower bound for the pseudo-Anosov dilatation unless
orientability of the invariant foliations is separately proved.

\subsection{The exact interface between the two braid layers}

There are two different fundamental groups in the example.  Moving the five
critical values gives $B_5$ and permutes the $1296$ LL sheets.  Holding $P$
fixed and moving $t$ in the complement of its critical values gives $F_5$ and
then a six-root braid.  There is no canonical homomorphism $B_5\to B_6$ joining
these two statements.

The correct common object is the mixed configuration space
\begin{equation}
 \mathcal U_{5,1}
 =\{(Q,t)\in\mathcal U_5\times\C:Q(t)\ne0\}.
 \label{eq:mixed-critical-value-configuration}
\end{equation}
Its fundamental group $B_{5,1}$ is the mixed braid group of five unordered
critical values and one distinguished moving parameter, and it fits into
\begin{equation}
 1\longrightarrow F_5\longrightarrow B_{5,1}
 \longrightarrow B_5\longrightarrow1.
 \label{eq:mixed-braid-exact-sequence}
\end{equation}
This is the Fadell--Neuwirth configuration-space exact sequence
\cite{FadellNeuwirth1962ConfigurationSpaces}.
Pull it back along $\operatorname{LL}_6$:
\begin{equation}
 \mathcal E_6^\circ
 =\{(P,t)\in\mathcal P_6^\circ\times\C:Q_P(t)\ne0\}.
 \label{eq:sextic-ll-moving-parameter-space}
\end{equation}
With compatible basepoints, put
\begin{equation}
 H_6=\pi_1(\mathcal P_6^\circ),
 \qquad
 G_6=\pi_1(\mathcal E_6^\circ).
 \label{eq:sextic-two-monodromy-groups}
\end{equation}
Then the covering pullback gives the exact diagram
\begin{equation}
\begin{array}{ccccccccc}
1&\longrightarrow&F_5&\longrightarrow&G_6&\longrightarrow&H_6&\longrightarrow&1\\
 &&\Vert&&\downarrow&&\downarrow&&\\
1&\longrightarrow&F_5&\longrightarrow&B_{5,1}&\longrightarrow&B_5&\longrightarrow&1,
\end{array}
\label{eq:sextic-two-braid-exact-diagram}
\end{equation}
where the two nonidentity vertical inclusions have index $1296$.  The root map
\begin{equation}
 \mathcal E_6^\circ\longrightarrow\UConf_6(\C),
 \qquad
 (P,t)\longmapsto P^{-1}(t),
 \label{eq:sextic-root-configuration-map}
\end{equation}
induces $G_6\to B_6$, whose restriction to the vertical $F_5$ over $P_0$ is
the surjection \eqref{eq:sextic-root-monodromy-surjection}.

Indeed, $H_6\subset B_5$ is the subgroup corresponding to the connected
$1296$-sheet LL cover, and
$G_6=(B_{5,1}\to B_5)^{-1}(H_6)$.  Surjectivity of
$B_{5,1}\to B_5$ gives
$[B_{5,1}:G_6]=[B_5:H_6]=1296$ and the exact top row.

Diagram \eqref{eq:sextic-two-braid-exact-diagram} is the precise
slice-to-slice relation in this example.  The lower row remembers only a moving
event schedule and a distinguished test value.  Its LL pullback remembers
which sextic realization is being transported.  The root map then extracts a
six-strand history.  A future AEG history functor would have to land in this
pullback or an equivalent groupoid; attaching a $B_6$ word after forgetting the
LL sheet would lose exactly the information at issue.

\subsection{An arithmetic fiber with no hidden endomorphisms}

The regular value $t=1$ gives the genus-two curve
\begin{equation}
 C_1:\quad y^2=x^6-x-1,
 \qquad
 J_1=\operatorname{Jac}(C_1).
 \label{eq:sextic-arithmetic-fiber}
\end{equation}

\begin{proposition}[Arithmetic maximality of one fiber]
\label{prop:sextic-arithmetic-maximality}
The arithmetic fiber satisfies
\begin{equation}
 \operatorname{Gal}(x^6-x-1/\Q)\cong S_6,
 \qquad
 \operatorname{End}_{\overline\Q}(J_1)=\Z.
 \label{eq:sextic-galois-endomorphism}
\end{equation}
Consequently $J_1$ is absolutely simple and has no complex multiplication.
The Galois action on $J_1[2]$ has full image
\begin{equation}
 S_6\cong\operatorname{Sp}_4(\F_2).
 \label{eq:sextic-arithmetic-two-torsion}
\end{equation}
This is the same finite target as the topological reduction in
\eqref{eq:sextic-mod-two-forgetting-sequence}, although the arithmetic Galois
action and a chosen topological loop are not thereby canonically identified
element by element.
\end{proposition}

\begin{proof}
Selmer proved the irreducibility of $x^n-x-1$ over $\Q$
\cite{Selmer1956IrreducibleTrinomials}.  Osada's trinomial criterion applies
to $n=6$ and gives the full symmetric Galois group
\cite{Osada1987GaloisTrinomials}; here its coprimality condition reduces to
$\gcd(5,6)=1$.  The discriminant is
\begin{equation}
 6^6+5^5=49781=67\cdot743,
 \label{eq:sextic-arithmetic-discriminant}
\end{equation}
so the fiber is smooth.  Zarhin's theorem for a hyperelliptic curve whose
defining polynomial has Galois group $S_n$ or $A_n$ gives the geometric
endomorphism ring $\Z$ in characteristic zero
\cite{Zarhin2000HyperellipticJacobians}.  The six Weierstrass points give the
standard faithful action of $S_6$ on the four-dimensional symplectic
$2$-torsion module, yielding \eqref{eq:sextic-arithmetic-two-torsion}.
\end{proof}

Independently, the integral Gauss--Manin image in
Theorem~\ref{thm:sextic-full-genus-two-monodromy} is maximal, and one
rational fiber has maximal mod-two arithmetic image and minimal possible
geometric endomorphism ring.  This does not identify arithmetic Frobenius with
topological monodromy without a comparison path and basepoint.  It says instead
that neither side is made rich by an accidental extra symmetry of the chosen
Jacobian.

\subsection{What the conjunction proves, and an explicit LL--Igusa twin}

The strongest proved conclusion is not that one has listed many familiar
objects.  It is the following conjunction for one sparse integral pencil:
\begin{equation}
 \boxed{
 \begin{gathered}
 Q_0(t)=6^{-6}\operatorname{disc}_x(x^6-x-t)
 \text{ gives all singular parameter values},\\
 \pi_1(\C\setminus V)\twoheadrightarrow B_6
 \twoheadrightarrow\operatorname{Mod}(\Sigma_2),\\
 \operatorname{im}\rho_{\mathrm{GM}}=\operatorname{Sp}_4(\Z),
 \qquad
 \operatorname{Gal}(x^6-x-1)=S_6,\\
 A_\gamma=\im(P_0-\gamma)
 \text{ and }K_\gamma=(P_0-\gamma)^{-1}(0)
 \text{ are one carrier--thread pair.}
 \end{gathered}}
 \label{eq:sextic-conjunction}
\end{equation}
It realizes the carrier--thread--spectral part of the original tube intuition:
one polynomial relation between slices supports all six-braids and all
genus-two gluings while its real part undergoes controlled saddle
reconnections.  It does not produce a proper zero tube, and it does not yet
validate arithmetic-history naturality.  The loop $\gamma$, or the LL sheet
$P$, has not yet been forced by a typed add--multiply history.

The finite LL fiber makes that missing naturality test concrete.  Its first
geometric part can in fact be completed without enumerating all $1296$ sheets.

Put
\begin{equation}
 R=-\frac{3^{12}5^{10}}{2^{14}},
 \qquad
 \alpha<0,
 \qquad
 \alpha^{15}=R,
 \qquad
 \delta\in\C,
 \qquad
 \delta^2=\frac{64\alpha^5}{253125},
 \label{eq:explicit-ll-twin-parameters}
\end{equation}
and choose either square root $\delta$.  Define
\begin{equation}
 P_1(x)=x^6+\alpha x^4+\frac4{15}\alpha^2x^2
       +\delta x+\frac8{675}\alpha^3.
 \label{eq:explicit-ll-twin-polynomial}
\end{equation}

\begin{theorem}[An explicit LL--Igusa twin]
\label{thm:explicit-ll-igusa-twin}
The two monic centered sextics $P_0=x^6-x$ and $P_1$ belong to distinct
$\mu_6$-orbits in the same reduced LL fiber:
\begin{equation}
 Q_{P_0}(t)=Q_{P_1}(t)=t^5+\frac{5^5}{6^6}.
 \label{eq:explicit-ll-twin-same-event-polynomial}
\end{equation}
Nevertheless the common regular slice $t=1$ gives nonisomorphic complex
genus-two curves:
\begin{equation}
 C_0:y^2=x^6-x-1,
 \qquad
 C_1:y^2=P_1(x)-1,
 \qquad
 C_0\not\cong_{\C}C_1.
 \label{eq:explicit-ll-twin-curves}
\end{equation}
Consequently the genus-two moduli reading of the slice does not factor through
the LL event polynomial, even on the single finite fiber over $Q_0$.
\end{theorem}

\begin{proof}
The sparse normalized polynomial
\begin{equation}
 S(x)=x^6+x^4+\frac4{15}x^2
      +\frac8{225\sqrt5}x+\frac8{675}
 \label{eq:explicit-ll-twin-normalized-polynomial}
\end{equation}
satisfies the exact Sylvester-resultant identity
\begin{equation}
 6^{-6}\operatorname{disc}_x(S(x)-t)
 =t^5-\frac{256}{3^{18}5^5}.
 \label{eq:explicit-ll-twin-normalized-resultant}
\end{equation}
One way to find it is to start with
$x^6+x^4+cx^2+dx+e$, set the four nonconstant lower coefficients of
its event quintic to zero, and eliminate: the reduced equations give
$15c-4=0$, $253125d^2-64=0$, and $e=8/675$.

For a scalar $\lambda\ne0$, the polynomial
$\lambda^6S(x/\lambda)$ has critical values $\lambda^6$ times those of
$S$.  Hence
\begin{equation}
 Q_{\lambda^6S(x/\lambda)}(t)
 =\lambda^{30}Q_S(t/\lambda^6)
 =t^5-\frac{256\lambda^{30}}{3^{18}5^5}.
 \label{eq:explicit-ll-twin-scaling}
\end{equation}
Choose $\lambda^2=\alpha$ and the sign of $\lambda$ so that
$\frac8{225\sqrt5}\lambda^5=\delta$.  The relation on $\delta^2$ makes
this possible, and $\lambda^{30}=\alpha^{15}=R$.  Equations
\eqref{eq:explicit-ll-twin-normalized-resultant}--
\eqref{eq:explicit-ll-twin-scaling} then give
\eqref{eq:explicit-ll-twin-same-event-polynomial}.  Since $Q_0$ is
square-free, $P_1$ lies in the finite \`etale locus and is a reduced LL sheet.
Its $x^4$ coefficient is nonzero, whereas every source rotation of $P_0$ has
zero $x^4$ coefficient, so the two sheets lie in distinct $\mu_6$-orbits.

It remains to compare the curves without choosing numerical roots.  For a
binary sextic
\begin{equation}
 F(X,Z)=\sum_{k=0}^6 c_kX^{6-k}Z^k
 \label{eq:binary-sextic-for-twin}
\end{equation}
set
\begin{equation}
 A(F)=\sum_{k=0}^6
       \frac{(-1)^kc_kc_{6-k}}{\binom6k},
 \qquad
 \mathcal J(F)=\frac{A(F)^5}{\operatorname{disc}(F)}.
 \label{eq:quadratic-clebsch-absolute-ratio}
\end{equation}
Up to one fixed nonzero normalization, $A$ is the quadratic
Igusa--Clebsch invariant $I_2$: explicitly,
$(F,F)_6=(6!)^2A(F)$.  Under a $\operatorname{GL}_2$ change of variables,
$A$ has determinant weight $6$ and the discriminant has determinant weight
$30$; under $F\mapsto rF$ they have degrees $2$ and $10$.  Thus
$\mathcal J$ is unchanged by both operations and is an absolute genus-two
isomorphism invariant \cite{Igusa1960ArithmeticVariety}.

Both sextics at $t=1$ have discriminant
\begin{equation}
 6^6Q_0(1)=6^6+5^5=49781,
 \label{eq:explicit-ll-twin-common-discriminant}
\end{equation}
whereas direct substitution in \eqref{eq:quadratic-clebsch-absolute-ratio}
gives
\begin{equation}
 A(P_0-1)=-2,
 \qquad
 A(P_1-1)=-2+\frac8{135}\alpha^3<-2.
 \label{eq:explicit-ll-twin-clebsch-values}
\end{equation}
The two values of $\mathcal J$ are therefore different.  Every isomorphism
between complex genus-two curves descends, by uniqueness of the hyperelliptic
involution, to a projective transformation of their branch lines.  The
$\mathcal J$ discrepancy consequently proves
\eqref{eq:explicit-ll-twin-curves}.  The exact resultant and invariant checks
are reproduced by
\texttt{paper-3/computations/ll-igusa-twin/verify-explicit-twin.py}.
\end{proof}

\subsection{What forgetting preserves: spectrum, variance, and divisor charge}

The explicit pair also gives a precise, limited sense in which information
forgotten by the LL map reappears as an energy.  Put
\begin{equation}
 \mathcal U_5^{(1)}
 =\{Q\in\mathcal U_5:Q(1)\ne0\},
 \qquad N=216,
 \label{eq:ll-energy-base}
\end{equation}
and quotient the degree-$1296$ LL cover by the free source-rotation action:
\begin{equation}
 \overline{\operatorname{LL}}_6:
 \overline{\mathcal P}_6^{(1)}
 =\operatorname{LL}_6^{-1}(\mathcal U_5^{(1)})/\mu_6
 \longrightarrow \mathcal U_5^{(1)}.
 \label{eq:ll-quotient-cover-for-energy}
\end{equation}
This is a finite \`etale cover of degree $N$.  For a point $[P]$ of its fiber,
let
\begin{equation}
 \jmath([P])
 =\mathcal J(P(x)-1)
 =\frac{A(P(x)-1)^5}{\operatorname{disc}(P(x)-1)}.
 \label{eq:ll-sheet-moduli-observable}
\end{equation}
The ratio is unchanged by $P(x)\mapsto P(\xi x)$, so it is well defined on
the quotient.

\begin{proposition}[Spectral descent and finite-fiber forgetting energy]
\label{prop:ll-forgetting-energy}
For $Q\in\mathcal U_5^{(1)}$, write
$\mathcal F_Q=\overline{\operatorname{LL}}_6^{-1}(Q)$ and set
\begin{equation}
 V_Q=\operatorname{Map}(\mathcal F_Q,\C),
 \qquad
 \operatorname{av}_Q(f)=\frac1N\sum_{p\in\mathcal F_Q}f(p),
 \qquad
 V_Q^0=\ker(\operatorname{av}_Q).
 \label{eq:ll-averaging-space}
\end{equation}
Then the permutation representation of LL monodromy preserves the canonically
split exact sequence
\begin{equation}
 0\longrightarrow V_Q^0\longrightarrow V_Q
 \xrightarrow{\ \operatorname{av}_Q\ }\C\longrightarrow0.
 \label{eq:ll-averaging-exact-sequence}
\end{equation}
The sheet observable $\jmath_Q=(\jmath(p))_{p\in\mathcal F_Q}$ has the
monodromy-invariant, nonnegative variance
\begin{align}
 \mathcal E_{\mathrm{LL}}(Q)
 &=\frac1N\sum_{p\in\mathcal F_Q}
   \left|\jmath(p)-\operatorname{av}_Q(\jmath_Q)\right|^2
 \label{eq:ll-forgetting-energy}\\
 &=\frac1{N^2}\sum_{\{p,q\}\subset\mathcal F_Q}
   |\jmath(p)-\jmath(q)|^2.
 \label{eq:ll-forgetting-energy-pairwise}
\end{align}
It vanishes exactly when this moduli reading is constant on the fiber.

The algebraic object retained after forgetting the sheet is the spectral
polynomial
\begin{equation}
 \Theta_Q(T)
 =\operatorname{Nm}_{\overline{\mathcal P}_6^{(1)}/\mathcal U_5^{(1)}}
   (T-\jmath)
 =\prod_{p\in\mathcal F_Q}(T-\jmath(p)).
 \label{eq:ll-moduli-spectral-polynomial}
\end{equation}
Its coefficients are regular single-valued functions of $Q$, given by traces,
norms, and higher symmetric functions of the sheet values.

At $Q_0$, Theorem~\ref{thm:explicit-ll-igusa-twin} gives the exact positive
lower bound
\begin{equation}
 \mathcal E_{\mathrm{LL}}(Q_0)
 \ge
 \frac{1}{216^2\,49781^2}
 \left|
  \left(-2+\frac8{135}\alpha^3\right)^5+32
 \right|^2
 >0.
 \label{eq:ll-forgetting-energy-lower-bound}
\end{equation}
In particular, $\Theta_{Q_0}$ is not a $216$th power of one linear factor.
\end{proposition}

\begin{proof}
The source-rotation action is free by
Theorem~\ref{thm:sextic-ll-event-schedule}, so
\eqref{eq:ll-quotient-cover-for-energy} is finite \`etale of degree $216$.
Moreover
\begin{equation}
 \operatorname{disc}(P(x)-1)=6^6Q(1),
 \label{eq:ll-energy-common-denominator}
\end{equation}
and the numerator in \eqref{eq:ll-sheet-moduli-observable} is polynomial in
the coefficients of $P$.  Thus $\jmath$ is regular on the stated open set.
Its invariance under source rotations follows from the
$\operatorname{GL}_2$ invariance proved in
Theorem~\ref{thm:explicit-ll-igusa-twin}.

Monodromy only permutes the $N$ points.  It therefore preserves averaging,
the constant functions, the zero-sum kernel, and the Hermitian product
\begin{equation}
 \langle f,g\rangle_Q
 =\frac1N\sum_{p\in\mathcal F_Q}f(p)\overline{g(p)}.
 \label{eq:ll-permutation-hermitian-product}
\end{equation}
Orthogonal projection onto the constants gives
\eqref{eq:ll-forgetting-energy}; the standard variance identity gives
\eqref{eq:ll-forgetting-energy-pairwise}.  Finite-flat norm descent gives
\eqref{eq:ll-moduli-spectral-polynomial}.

Finally, the two distinct source-rotation orbits already exhibited have values
\begin{equation}
 \jmath([P_0])=-\frac{32}{49781},
 \qquad
 \jmath([P_1])
 =\frac{\left(-2+\frac8{135}\alpha^3\right)^5}{49781}.
 \label{eq:ll-twin-sheet-observable-values}
\end{equation}
Their pair is one term of the nonnegative sum in
\eqref{eq:ll-forgetting-energy-pairwise}, proving
\eqref{eq:ll-forgetting-energy-lower-bound}.
\end{proof}

The word \emph{energy} in Proposition~\ref{prop:ll-forgetting-energy} has a
strict boundary.  The function $\mathcal E_{\mathrm{LL}}$ is independent of
sheet labels and returns to the same value around a closed LL loop, but it need
not be constant along an arbitrary open path in $\mathcal U_5^{(1)}$ because
$\jmath_Q-\operatorname{av}_Q(\jmath_Q)\mathbf1$ is not generally a flat
section.  On genus-two homology, the genuinely Gauss--Manin-conserved form is
the alternating intersection pairing.  Full
$\operatorname{Sp}_4(\Z)$ monodromy precludes a positive-definite form fixed by
all monodromy.  Thus \eqref{eq:ll-forgetting-energy} is a finite-fiber
information variance, not a Noether or Dirichlet energy.

There is nevertheless an exact conserved charge along the two displayed
$t$-pencils.  Let
\begin{equation}
 \mathcal J_i(t)=\mathcal J(P_i(x)-t),
 \qquad
 \beta=\frac4{135}\alpha^3<0.
 \label{eq:ll-twin-moving-invariants}
\end{equation}

\begin{corollary}[Displaced invariant divisor]
\label{cor:ll-twin-displaced-invariant-divisor}
The two pencils satisfy
\begin{equation}
 A(P_0-t)=-2t,
 \qquad
 A(P_1-t)=-2(t-\beta),
 \qquad
 \frac{\mathcal J_1(t)}{\mathcal J_0(t)}
 =\left(\frac{t-\beta}{t}\right)^5.
 \label{eq:ll-twin-moving-invariant-ratio}
\end{equation}
Both $t=0$ and $t=\beta$ are regular fibers, and on $\P^1_t$
\begin{equation}
 \operatorname{div}\!\left(\frac{\mathcal J_1}{\mathcal J_0}\right)
 =5[\beta]-5[0].
 \label{eq:ll-twin-divisor-charge}
\end{equation}
Equivalently, the logarithmic derivative has residues $+5$ and $-5$ and total
residue zero.
\end{corollary}

\begin{proof}
Direct substitution into \eqref{eq:quadratic-clebsch-absolute-ratio} gives the
first two formulas.  The discriminants cancel because both event polynomials
are $Q_0$.  Also $Q_0(0)\ne0$, while
\begin{equation}
 \beta^5
 =-\frac{5^5}{2^4 3^3}
 \ne-\frac{5^5}{6^6},
 \label{eq:ll-twin-beta-not-discriminant}
\end{equation}
so $Q_0(\beta)\ne0$.  The divisor and residue statements follow from the
last expression in \eqref{eq:ll-twin-moving-invariant-ratio}.
\end{proof}

Thus one and the same forgotten difference has two complementary forms.  Across
the finite LL fiber it is the zero-sum component in
\eqref{eq:ll-averaging-exact-sequence}, with positive norm
\eqref{eq:ll-forgetting-energy}; along the slice parameter it is the displacement
of a smooth internal invariant divisor, with the balanced charge
\eqref{eq:ll-twin-divisor-charge}.  A Hodge- or Siegel-metric energy comparison
would be a deeper analytic refinement and remains part of
Problem~\ref{prob:ll-igusa-twin-test}.

\begin{remark}[Relation to recent work of Farb and collaborators]
\label{rem:sextic-farb-relation}
Farb--Kisin--Wolfson relate the general sextic to level-$2$
hyperelliptic modular functions and use the exceptional
$S_6\cong\operatorname{Sp}_4(\F_2)$ representation in resolvent problems
\cite{FarbKisinWolfson2023Modular}.  Farb--Wolfson study the resolvent degree
of the same sextic/configuration-space geometry
\cite{FarbWolfson2019Resolvent}, while Farb--Looijenga's Wiman--Edge work is a
model for reading geometry, monodromy, arithmetic, and periods from one pencil
\cite{FarbLooijenga2021WimanEdge}.  By conceptual analogy, Farb's rigidity
program supplies a broader language for asking when a moduli reading might
factor through a coarser algebro-geometric construction
\cite{Farb2024ModuliRigidity}.
Their more recent relative-essential-dimension framework asks how much finite
branched auxiliary data can compress a covering problem
\cite{FarbWolfson2025RelativeEssentialDimension}; applying that question to the
LL forgetting cover is a possible next step, not an input used below.

Those works cover much of the lower passage from six branch points to
genus-two modular and symplectic data.  They do not compute the LL fiber used
here.  The statement established in
Theorem~\ref{thm:explicit-ll-igusa-twin} is the
upper, fiberwise nonfactorization: one fixed critical-value event polynomial
has two explicit LL realizations with distinct genus-two moduli readings.  The
Igusa quartic occurring in the level-$2$ Satake compactification and the
Igusa--Clebsch invariant used in \eqref{eq:quadratic-clebsch-absolute-ratio}
belong to the same genus-two invariant geometry but are not the same object.
\end{remark}

\begin{problem}[The full LL--Igusa census and deeper twins]
\label{prob:ll-igusa-twin-test}
Enumerate the $1296$ monic centered sheets over $Q_0$, quotient the explicit
$\mu_6$ source rotations, and determine the full image of the resulting $216$
orbits in $\mathcal M_2$ at $t=1$.  Compare their complete
weighted-projective Igusa--Clebsch points
$[I_2:I_4:I_6:I_{10}]$; only on a chart where the chosen denominators are
nonzero should these be replaced by absolute affine invariants.

Theorem~\ref{thm:explicit-ll-igusa-twin} settles existence of distinct
geometric moduli points.  Still open is the finer census: place candidate
pairs in a common number field, compare Frobenius characteristic polynomials
at the same good primes, and compare reduced Siegel period matrices after a
complex embedding.  A stronger marked test fixes a common critical-value
ordering, meridian system, and homology marking, groups sheets by simultaneous
conjugacy of their local $\operatorname{Sp}_4(\Z)$ transvections, and asks
whether distinct Igusa points survive inside one such group after quotienting
the Hurwitz action.  This would test whether even
\begin{equation}
 Q_P+\text{coarse homological monodromy}
 \label{eq:ll-igusa-strong-forgetting-test}
\end{equation}
fails to recover the geometry forgotten by $P\mapsto Q_P$.
\end{problem}

The explicit twin is a genuine finite, exact instance of information loss:
the event schedule does not determine the slice's complex geometry.  It does
not yet make the sextic or its LL sheet natural in unrestricted AEG history,
and it does not turn the remaining arithmetic and period comparisons into
automatic consequences.
